% page 397 of the facsimile (image p-396 of the scan)
% block 167.6 x 246.3 mm ; left margin 34.4 mm
\pgc{12.700mm}{0.000mm}{
\l{\marge[78.511mm]{22.296mm}{\ornamento{11.819mm}{ornaments/letters/letter-N.png}}\cel{52}389}
\l{}
\l{}
\l{}
\l{}
\l{}
\l{}
\l{}
\l{}
\l{-n. (\sou{gram.)}\cel{3}Signo qua indikas ke o la komplemento direta o}
\l{\cel{3}la atributo esas inversigita en la propoziciono - e, per lo,}
\l{\cel{3}preventas kompren-erori che la lektanto o la askoltanto.}
\l{}
\l{nabo. Centro di roto, adube konvergas la radii e quan trairas}
\l{\cel{3}la rotaxo.- DE.}
\l{}
\l{nabeto.\cel{3}(\sou{bot.)}\cel{2}Varietato di napo, di qua la grani esas oleoza.}
\l{\cel{3}L.\cel{1}\sou{brassica rapa oleifera}. - FS.}
\l{}
\l{\sou{nabobo}.\cel{2}I. Princo mohamedista en India. - II. Personego qua re-}
\l{\cel{3}trovenis de India (o de altra landi exotika) kun havajo grand-}
\l{\cel{3}ega. - DEFIRS.}
\l{}
\l{\sou{naciono}. I. Asemblitaro de homi ye qua konsistas socio politi-}
\l{\cel{3}kala, qua esas guvernata, administrata o jerata da institu-}
\l{\cel{3}curi komuna. -II. Singla ek la dividuri en qui repartisesis}
\l{\cel{3}olim la studenti di la universitato di la urbo Paris. - DEFIRS.}
\l{}
\l{\sou{nadiro}. (\sou{astron.}) Punto di la cielo-sfero ube abutas vertikalo}
\l{\cel{3}qua departas de la loko ye qua onu esas, e pasante tra la}
\l{\cel{3}centro di la Terglobo. - DEFIRS.}
\l{}
\l{\sou{nafto}. (\sou{kemio}).Hidrokarbido, sorto di bitumo liquida, volatila}
\l{\cel{3}ed inflamebla. - DEFIRS.}
\l{}
\l{\sou{naftalino}\cel{4}(\sou{kemio)}. Hidrokarbido quan produktas la distilo di}
\l{\cel{3}minkarbono, di gudro. - DEFIRS.}
\l{}
\l{\sou{naiva}. Qua esas tam kredera kam mento sen-experienca. - DEFR.}
\l{}
\l{\sou{nakarato.}\cel{2}Koloro reda oranjea. - DEFS.}
\l{}
\l{\sou{naktigalo}. (\sou{zool.)}\cel{3}Mikra ucelo, ek la klaso \textquotedbl{}paseri\textquotedbl{}, di qua}
\l{\cel{3}la plumaro esas griza, famoza pro la beleso di lua kanto-}
\l{\cel{3}stilo. - L.\cel{1}\sou{sylvia luscinia.}\cel{2}- DE.}
\l{}
\l{\sou{nam}.\cel{2}Konjunciono qua unionas, a propoziciono, propoziciono qua}
\l{\cel{3}sequas olca e qua donas la \textquotedbl{}pro quo\textquotedbl{} di to quon afirmas la}
\l{\cel{3}unesma. - L.}
\l{}
\l{\sou{nano}. Persono qua esas mikra extraordinare. - FIS.}
\l{}
\l{\sou{nankino}.\cel{2}Telo de kotono, di qua la koloro esas uniforma, maxim-}
\l{\cel{3}multa-kaze flava klare. - DEFIR.}
\l{}
\l{\sou{napo}. (\sou{bot.})\cel{2}Planto krucifera, kun radiko pulpoza e nutriva.}
\l{\cel{3}- L.\cel{1}\sou{brassica napus esculenta}.\cel{2}FIS.}
}
